% Mục 4.1: Chiến lược kiểm tra. Nguồn: tests/*.cpp, mk/part{1,2,3}.mk,
% results/{apple-m5,C4M7Y670NT,LAP1184}/{test,tsan,mutation}.log.

Nhóm kiểm hai nhóm tính chất. Tính an toàn (safety) nghĩa là không có hai thread
cùng ở trong critical section, không thread nào ra khỏi barrier sớm, không giá trị counter nào bị trả
về hai lần. Tính sống còn (liveness) nghĩa là chương trình không treo. Công bằng chỉ được đo.

Cùng một bài kiểm tra dạng template chạy cho mọi cài đặt, kể cả lock và barrier của thư viện.

Nhóm dùng bốn kỹ thuật:
\begin{itemize}[nosep]
  \item stress test với ít thread (1--4), với số luồng phần cứng và với gấp đôi số đó (\PCname{apple-m5}:
    10 và 20; \PCname{C4M7Y670NT} và \PCname{LAP1184}: 12 và 24); ở mức gấp đôi, thread bị tước core giữa giao thức;
  \item watchdog, thread giám sát báo \texttt{FAIL} khi bài kiểm tra ngừng tiến triển;
  \item ThreadSanitizer (TSan), công cụ của compiler báo hai truy cập xung đột, ít nhất một không atomic,
    mà không có quan hệ happens-before (atomic riêng của bài kiểm tra dùng \texttt{relaxed} để không tự tạo
    thứ tự, trừ đồng hồ \texttt{seq\_cst} của lượt kiểm thời gian thực);
  \item mutation testing, chạy bộ kiểm tra trên các biến thể cố ý làm hỏng.
\end{itemize}
Barrier dùng thêm schedule fuzzing. Counter dùng thêm phép kiểm thứ tự thời gian thực.

Bảng~\ref{tab:corr:summary} (phụ lục) tóm tắt kết quả trên ba máy. Mọi cài đặt đạt mọi bài kiểm tra, TSan
không báo data race nào trên mã đúng, và bộ kiểm tra bắt cả 14 biến thể hỏng. Bản TSan chạy với kích thước nhỏ hơn
(ít vòng hơn 10--20 lần) vì TSan làm chương trình chậm đi nhiều. Đầu ra nằm trong
\texttt{results/<máy>/}.
% nguồn: tests/test_locks.cpp (small: chia 10, 100 ms), tests/test_counter.cpp (small: chia 16),
% tests/test_barrier.cpp (small: g_size 0,05),
% mk/*.mk (make tsan chạy với đối số small).

